\section{The earlier accuracy lower bound}\label{app:accuracy}

This appendix supplies the $k/\eps^2$ argument used in the main theorem.
It retains the definitions and constants of the earlier note.

\subsection{Information in adaptive shifted-GOE queries}

For a symmetric matrix $G$, use the GOE normalization
\begin{equation}\label{eq:goe}
 G_{ij}=G_{ji}\sim N(0,1/n)\quad(i<j),\qquad
 G_{ii}\sim N(0,2/n),
\end{equation}
with independent upper-triangular entries. Its density on the vector space of
symmetric matrices is proportional to
$\exp(-n\norm{G}_{\F}^2/4)$. In particular, it is isotropic for the Frobenius
inner product on that vector space.

\begin{lemma}[Adaptive transcript divergence]\label{lem:kl}
Fix symmetric $A_0,S\in\R^{n\times n}$ with $\norm{S}_{\op}\le1$ and a real
$\delta$. Let $\mathbb P_S$ be the law of the transcript of any randomized
adaptive algorithm using at most $q\le n$ vector queries on
\[
 A=A_0+\delta S+G.
\]
Let $\mathbb P_0$ be its transcript law on $A_0+G$. Include the independent
random seed in the transcript. Then
\begin{equation}\label{eq:kl}
 \KL(\mathbb P_S\Vert\mathbb P_0)\le\frac{n\delta^2q}{2}.
\end{equation}
The statement permits queries on either side of $A$.
\end{lemma}

\begin{proof}
\emph{Reduction to orthonormal fresh directions.}
All matrices in question are symmetric, so choosing the transpose side gives
no additional type of measurement. Subtract the known product $A_0v$ from
each answer. If a new query vector has a component in the span of previous
queries, its image on that component is already known. Removing it and
normalizing the remaining vector recovers the same fresh information using
one call. Queries wholly in the old span give no information and can be
omitted. We may therefore work with an orthonormal sequence of fresh vectors,
chosen adaptively. An algorithm that stops early can be padded with
orthogonal queries. Data processing then bounds its original transcript by
the padded transcript. Fix the seed for the moment.

\emph{Gaussian exposure under adaptivity.}
After $t-1$ fresh directions, let $V$ be the matrix of those directions and
$\Pi=VV^{\mathsf T}$. For a fixed value of $S$, the observed answers determine
$Z=GV$. The symmetric component of $G$ determined by these observations is
\begin{equation}\label{eq:observedcomponent}
 G_{\rm obs}=ZV^{\mathsf T}+VZ^{\mathsf T}
             -V(V^{\mathsf T}Z)V^{\mathsf T}.
\end{equation}
The undetermined part is $(I-\Pi)G(I-\Pi)$. These are orthogonal subspaces of
the space of symmetric matrices for the Frobenius inner product. Conditional
on the history, the latter part remains a centered GOE compression on
$\operatorname{range}(I-\Pi)$, with the normalization $1/n$, not the inverse
of the remaining dimension.

For completeness, this claim also holds when $V$ is adaptive. Initially it
is just the orthogonal decomposition of an isotropic Gaussian. Given the
past, the next direction is fixed. In an orthonormal basis of the remaining
space beginning with that direction, the new Gaussian column and the
unexposed lower-right symmetric block are independent. Conditioning on the
column leaves that block unchanged. Induction proves the assertion. The
choice of the next vector is a function of already conditioned-on answers,
so it does not impose an additional constraint on the unexplored block.

\emph{One-step divergence.}
Let $v$ be the next unit vector, with $\Pi v=0$. The new information is the
component
\[
 y=(I-\Pi)(A-A_0)v.
\]
The component in $\operatorname{range}(\Pi)$ is already determined by the old
answers and symmetry:
$V^{\mathsf T}(A-A_0)v=((A-A_0)V)^{\mathsf T}v$.
At a fixed history, the conditional mean of $y$ under the shifted law is
$\delta(I-\Pi)Sv$, and under the baseline it is zero. In both cases its
conditional covariance, restricted to $\operatorname{range}(I-\Pi)$, is
\begin{equation}\label{eq:covariance}
 \Sigma_v=\frac1n\bigl((I-\Pi)+vv^{\mathsf T}\bigr).
\end{equation}
The variance is $2/n$ in direction $v$ and $1/n$ in the orthogonal remaining
directions. In particular, $\Sigma_v^{-1}\preceq nI$ on its support.

Two Gaussians with identical covariance and mean difference $m$ have
divergence $m^{\mathsf T}\Sigma_v^{-1}m/2$, as follows directly by integrating
the logarithm of their density ratio. The conditional divergence is therefore
at most
\[
 \frac{n\delta^2}{2}\norm{(I-\Pi)Sv}^2
 \le\frac{n\delta^2}{2}.
\]
Applying the chain rule for relative entropy to the successive fresh answers
and summing gives \eqref{eq:kl}. Integrating over the seed gives the same bound
for randomized algorithms, since its distribution is identical in both
experiments.
\end{proof}

\begin{remark}[Why exact arithmetic does not invalidate the bound]
This is a divergence calculation for the laws of exact real-valued answers.
No rounding, finite-precision hypothesis, bounded query coefficients, or
nonadaptive restriction is used. Normalizing a query loses no information
because its scale is known. The unexplored Gaussian compression is what
prevents arbitrary measurable postprocessing from revealing the hidden
parameter for free.
\end{remark}

\subsection{A local approximation-to-subspace inequality}

The next lemma is the deterministic part of the construction. It handles
arbitrary perturbations, nonsymmetric estimates, and estimates of rank
strictly smaller than $k$.

\begin{lemma}[Local loss]\label{lem:local}
Let $P$ be a rank-$k$ orthogonal projector on $\R^{2k}$, and let
$0<\delta\le1/2$. Put
\[
 C=(1-\delta)I+2\delta P.
\]
For any $W\in\R^{2k\times2k}$ and any matrix $D$ of rank at most $k$, choose
a rank-$k$ orthogonal projector $Q$ whose range contains the row space of $D$.
Then
\begin{align}\label{eq:local}
 &\norm{C+W-D}_{\F}^2-\tau_k(C+W)^2\notag\\
 &\hspace{1em}\ge
 \delta\norm{P-Q}_{\F}^2
 -\left(\frac{(1+\delta)^2}{\delta}+1\right)\norm{W}_{\F}^2.
\end{align}
\end{lemma}

\begin{proof}
Write $d=\norm{P-Q}_{\F}$ and $w=\norm{W}_{\F}$. Since $D=DQ$,
\[
 \norm{C+W-D}_{\F}^2\ge\norm{(C+W)(I-Q)}_{\F}^2.
\]
The rank-$k$ matrix $(1+\delta)P=CP$ is a valid competitor, so
\[
 \tau_k(C+W)^2\le\norm{C(I-P)+W}_{\F}^2.
\]
Subtracting these inequalities and expanding squared norms yields
\begin{align}\label{eq:localexpand}
 &\norm{C(I-Q)}_{\F}^2-\norm{C(I-P)}_{\F}^2
  +2\ip{W}{C(P-Q)}_{\F}-\norm{WQ}_{\F}^2.
\end{align}
Because $C^2=(1-\delta)^2I+4\delta P$, the clean difference is
\[
 \tr((P-Q)C^2)=4\delta(k-\tr(PQ))=2\delta d^2.
\]
The cross term is bounded below by $-2(1+\delta)wd$, and the last term is at
least $-w^2$. Thus \eqref{eq:localexpand} is bounded below by
\[
 2\delta d^2-2(1+\delta)wd-w^2.
\]
Use
$2(1+\delta)wd\le\delta d^2+(1+\delta)^2w^2/\delta$
to obtain \eqref{eq:local}.
\end{proof}

The projector $Q$ can be selected measurably from $D$: choose a deterministic
orthonormal basis of its row space, then extend it by a fixed coordinate-based
Gram--Schmidt rule until its dimension is $k$. No optimality, symmetry, or
uniqueness assumption on $D$ is needed.

\subsection{An explicit finite packing}

All logarithms in this section and in the information calculation are natural.
The metric on tuples of projectors is the square root of the sum of their
squared Frobenius distances.

\begin{lemma}[Projector tuples]\label{lem:packing}
For every pair of positive integers $k,M$, there is a finite collection
\[
 \bigl\{(P_1^\theta,\ldots,P_M^\theta):\theta=1,\ldots,N\bigr\}
\]
of rank-$k$ orthogonal projectors on $\R^{2k}$ such that
\begin{equation}\label{eq:packsize}
 \log N\ge\frac{Mk^2}{200}
\end{equation}
and, whenever $\theta\ne\theta'$,
\begin{equation}\label{eq:packdist}
 \sum_{j=1}^M\norm{P_j^\theta-P_j^{\theta'}}_{\F}^2
 \ge\frac{Mk}{1600}.
\end{equation}
\end{lemma}

\begin{proof}
\emph{Step 1: a single-block alphabet.}
Suppose first that $k\ge4$. Let $S$ range over sign matrices with independent
entries $\pm1/\sqrt{k}$. For fixed unit vectors $x,y$, the random variable
$x^{\mathsf T}Sy$ has moment-generating function at most
$\exp(t^2/(2k))$, because $\cosh u\le\exp(u^2/2)$ and
$\sum_{ij}x_i^2y_j^2=1$. Thus
\[
 \Prob\{|x^{\mathsf T}Sy|>a\}\le2\exp(-ka^2/2).
\]
A $1/4$-net of the unit sphere has size at most $9^k$: take a maximal separated
set and compare disjoint balls of radius $1/8$ with the enclosing ball of
radius $9/8$. For nets in both variables,
$\norm{S}_{\op}\le2\max_{x,y\text{ in nets}}|x^{\mathsf T}Sy|$.
Consequently,
\[
 \Prob\{\norm{S}_{\op}>10\}
 \le2\exp\bigl(-(25/2-2\log9)k\bigr)<1/2.
\]
At least $2^{k^2-1}$ sign matrices are therefore good.

Let $h(p)=-p\log p-(1-p)\log(1-p)$. A binary Hamming ball of radius
$\lfloor k^2/4\rfloor$ has size at most $\exp(h(1/4)k^2)$. To verify the usual
binomial bound, multiply each summand by
$p^i(1-p)^{d-i}$ with $d=k^2$ and $p=1/4$; for $i\le pd$ this weight is at
least $\exp(-dh(p))$, while the weighted sum is at most one.
Greedy deletion of such balls from the good matrices gives a set with mutual
Hamming distance greater than $k^2/4$ and log cardinality at least
\[
 (\log2-h(1/4))k^2-\log2\ge k^2/20,\qquad k\ge4.
\]

For each retained sign matrix set $T=S/20$, so $\norm{T}_{\op}\le1/2$, and let
$P_T$ project onto the graph of $T$, the column space of
$[I\;T^{\mathsf T}]^{\mathsf T}$. For two matrices $T,U$, an orthonormal basis
of the graph of $U$ is
$[I\;U^{\mathsf T}]^{\mathsf T}(I+U^{\mathsf T}U)^{-1/2}$; an orthonormal
basis of the orthogonal complement of the graph of $T$ is
$[-T\;I]^{\mathsf T}(I+TT^{\mathsf T})^{-1/2}$. Hence
\begin{align}\label{eq:graphidentity}
 \frac12\norm{P_T-P_U}_{\F}^2
 =\norm{(I+TT^{\mathsf T})^{-1/2}(U-T)
          (I+U^{\mathsf T}U)^{-1/2}}_{\F}^2.
\end{align}
Both inverse square roots have minimum singular value at least
$1/\sqrt{1+1/4}$. The Hamming separation also gives
$\norm{T-U}_{\F}^2>k/400$. It follows that
\[
 \norm{P_T-P_U}_{\F}^2
 \ge\frac{2}{(1+1/4)^2}\norm{T-U}_{\F}^2\ge k/400.
\]
For $k=1,2,3$, instead take the two complementary coordinate projectors.
Their squared distance is $2k$ and their log cardinality is $\log2\ge k^2/20$.
Thus for every $k$ there is a single-block alphabet of size $Q\ge2$ with
$\log Q\ge k^2/20$ and squared separation at least $k/400$.

\emph{Step 2: a product code.}
In the $Q$-ary cube of length $M$, a Hamming ball of radius $\lfloor M/4\rfloor$
has size at most
\[
 \exp\bigl(Mh(1/4)+(M/4)\log(Q-1)\bigr).
\]
Indeed, factor $(Q-1)^{M/4}$ out of the relevant binomial sum and apply the
binary bound just proved. Greedy packing then gives mutual Hamming distance
greater than $M/4$ and
\[
 \log N\ge M\bigl(\log Q-h(1/4)-\tfrac14\log(Q-1)\bigr)
 \ge\tfrac1{10}M\log Q\ge Mk^2/200.
\]
The middle inequality holds for $Q\ge2$: the function
$0.9\log Q-0.25\log(Q-1)$ is increasing there, and its value at $2$ is greater
than $h(1/4)$. Each differing coordinate contributes at least $k/400$ to
squared projector distance. This proves \eqref{eq:packdist}.
\end{proof}

\subsection{The hard distribution and the accuracy lower bound}

\subsubsection{Matrices indexed by the packing}

Set $M=n/(4k)=2^{L-2}$. Partition the coordinates into $M$ consecutive nodes
of size $4k$, and split each into two sets $I_j,J_j$ of size $2k$.
These are actual sibling blocks of the HODLR partition. Crucially, they are
not the unrestricted $k\times k$ leaf blocks.

Take the packing from Lemma~\ref{lem:packing}, and put
\begin{align}\label{eq:hardfamily}
 R_j^\theta&=2P_j^\theta-I_{2k},\notag\\
 A_0&=\operatorname{blockdiag}_{j=1}^M
       \begin{bmatrix}0&I_{2k}\\ I_{2k}&0\end{bmatrix},\notag\\
 S_\theta&=\operatorname{blockdiag}_{j=1}^M
       \begin{bmatrix}0&R_j^\theta\\ R_j^\theta&0\end{bmatrix},\qquad
 A=A_0+\delta S_\theta+G.
\end{align}
Each $R_j^\theta$ is a symmetric orthogonal matrix. Thus $S_\theta$ is symmetric
orthogonal too, and $\norm{S_\theta}_{\op}=1$.
The upper targeted block of the clean matrix is
\[
 C_j=(1-\delta)I_{2k}+2\delta P_j^\theta.
\]
Its singular values are $1+\delta$ and $1-\delta$, each repeated $k$ times.
The lower targeted block is its transpose. All other clean blocks are zero.
By Lemma~\ref{lem:separable},
\begin{equation}\label{eq:cleanopt}
 \OPT(A_0+\delta S_\theta)^2
 =2Mk(1-\delta)^2=\frac n2(1-\delta)^2.
\end{equation}
The packing parameters become
\begin{equation}\label{eq:packingn}
 \log N\ge\frac{nk}{800},\qquad
 \sum_j\norm{P_j^\theta-P_j^{\theta'}}_{\F}^2\ge\frac n{6400}.
\end{equation}

\subsubsection{A good-noise event}

Define
\[
 Z=\norm{G}_{\F}^2,\qquad
 T=\sum_{j=1}^M\norm{G_{I_j,J_j}}_{\F}^2.
\]
The GOE normalization gives
\[
 \E Z=n+1\le\tfrac54n,\qquad \E T=k.
\]
For the second identity, there are $M(2k)^2=nk$ entries in the targeted
\emph{upper} blocks, each with variance $1/n$. These entries are disjoint.
Markov's inequality and a union bound imply that the event
\begin{equation}\label{eq:event}
 \mathcal E=\{Z\le10n,\ T\le10k\}
\end{equation}
has probability at least
\begin{equation}\label{eq:eventprob}
 \Prob(\mathcal E)\ge1-\frac18-\frac1{10}=\frac{31}{40}.
\end{equation}
On this event, use the best clean HODLR approximation as a competitor for the
noisy matrix and apply $\norm{X+Y}_{\F}^2\le2\norm X_{\F}^2+2\norm Y_{\F}^2$:
\begin{equation}\label{eq:noisyopt}
 \OPT(A)^2\le n(1-\delta)^2+2Z\le21n.
\end{equation}
No independence between $Z$ and $T$ is needed.

\subsubsection{Decoding the hidden projectors}

Suppose an approximation algorithm succeeds on an input from
\eqref{eq:hardfamily}. Choose $Q_j$ to be a rank-$k$ projector containing the
row space of $B_{I_j,J_j}$, using the measurable rule following
Lemma~\ref{lem:local}. Let
\[
 r^2=\sum_{j=1}^M\norm{P_j^\theta-Q_j}_{\F}^2.
\]
Since $\eps<1/2$, success and \eqref{eq:noisyopt} imply, on $\mathcal E$,
\begin{align}\label{eq:globalexcess}
 \norm{A-B}_{\F}^2-\OPT(A)^2
 &\le (2\eps+\eps^2)\OPT(A)^2\notag\\
 &\le\frac{105}{2}\eps n.
\end{align}
Sum Lemma~\ref{lem:local} over the targeted upper blocks. Their excesses are
nonnegative and sum to at most the global excess by
Lemma~\ref{lem:separable}. Also,
\[
 \frac{(1+\delta)^2}{\delta}+1
 =\frac{1+3\delta+\delta^2}{\delta}\le\frac3\delta,
 \qquad 0<\delta\le1/2.
\]
It follows that
\begin{equation}\label{eq:radius}
 r^2\le\frac{105\eps n}{2\delta}+\frac{30k}{\delta^2}.
\end{equation}

Now set
\begin{equation}\label{eq:constants}
 D=2^{24},\qquad \eps_0=\frac1{2D},\qquad \delta=D\eps,
\end{equation}
and first restrict to
\begin{equation}\label{eq:coregime}
 0<\eps\le\eps_0,\qquad n\ge k/\eps^2.
\end{equation}
Then $0<\delta\le1/2$ and
\begin{equation}\label{eq:smallradius}
 \frac{r^2}{n}\le\frac{105}{2D}+\frac{30}{D^2}
 <\frac1{102400}<\frac1{25600}.
\end{equation}
The smallest squared separation in \eqref{eq:packingn} is at least $n/6400$.
Thus $r$ is less than half the minimum separation. Nearest-neighbor decoding
of the tuple $(Q_1,\ldots,Q_M)$ in the finite packing recovers $\theta$
uniquely. This is a measurable operation, and no query is needed to do it.

Let $\theta$ be uniform on the packing, independently of $G$ and the algorithm's
seed. The fixed-input guarantee \eqref{eq:guarantee} implies average success
at least $0.99$ under this distribution. The success event and $\mathcal E$
therefore intersect with probability at least
\begin{equation}\label{eq:decodeprob}
 \frac{31}{40}-\frac1{100}=\frac{153}{200}.
\end{equation}
Hence the decoding error probability is at most $47/200$.

\subsubsection{Information cannot support that decoding probability}

Let $\mathcal T$ denote the transcript including the seed. By
Lemma~\ref{lem:kl}, and using the baseline transcript law generated by $A_0+G$,
\begin{equation}\label{eq:mutualinfo}
 I(\theta;\mathcal T)
 \le\frac1N\sum_{\theta=1}^N
       \KL(\mathbb P_\theta\Vert\mathbb P_0)
 \le\frac{n\delta^2q}{2}.
\end{equation}
The first inequality follows from the identity
\[
 \frac1N\sum_\theta\KL(\mathbb P_\theta\Vert\mathbb P_0)
 =I(\theta;\mathcal T)+\KL(\overline{\mathbb P}\Vert\mathbb P_0),
\]
where $\overline{\mathbb P}$ is the mixture transcript law.

We recall the needed entropy bound directly. If $\widehat\theta$ is any
transcript-based decoder and $p_e=\Prob(\widehat\theta\ne\theta)$, conditioning
also on the error indicator gives
\[
 H(\theta\mid\mathcal T)\le\log2+p_e\log N.
\]
Since $H(\theta)=\log N$, this yields Fano's inequality in the form
\begin{equation}\label{eq:fano}
 p_e\ge1-\frac{I(\theta;\mathcal T)+\log2}{\log N}.
\end{equation}

Assume, for a contradiction, that
\begin{equation}\label{eq:qthreshold}
 q\le\frac{k}{6400\delta^2}.
\end{equation}
This bound is less than $n$ under \eqref{eq:coregime}, so the hypothesis
$q\le n$ in Lemma~\ref{lem:kl} applies. Equations \eqref{eq:mutualinfo} and
\eqref{eq:packingn} give $I(\theta;\mathcal T)\le(\log N)/16$.
Moreover, \eqref{eq:coregime} implies
$nk\ge4D^2$, so $\log N\ge nk/800\ge16\log2$.
Equation \eqref{eq:fano} therefore gives $p_e\ge7/8$, contradicting
$p_e\le47/200$ from \eqref{eq:decodeprob}. We have proved the following.

\begin{proposition}[Accuracy lower bound in the core regime]\label{prop:core}
If $\eps\le\eps_0$ and $n\ge k/\eps^2$, then
\begin{equation}\label{eq:corelower}
 q_*(n,k,\eps)\ge c_*\frac{k}{\eps^2},\qquad c_*=(6400D^2)^{-1}.
\end{equation}
\end{proposition}


\subsection{Extending the accuracy lower bound to all parameters}

Recall $D=2^{24}$, $\eps_0=1/(2D)$, and $c_*=(6400D^2)^{-1}$.
The function $q_*(n,k,\eps)$ is nonincreasing in $\eps$.
For $\eps\ge\eps_0$,
\[
 c_*\min\{n,k/\eps^2\}\le(c_*/\eps_0^2)k=k/1600,
\]
so Proposition~\ref{prop:exactlower} suffices.
For $\eps<\eps_0$ and $n\ge k/\eps^2$, apply Proposition~\ref{prop:core}.
Otherwise put $\zeta=\sqrt{k/n}>\eps$. If $\zeta\le\eps_0$, monotonicity
and the core bound at accuracy $\zeta$ give $q_*\ge c_*k/\zeta^2=c_*n$.
If $\zeta>\eps_0$, then $n<k/\eps_0^2$, and the exact-recovery obstruction
again gives $q_*\ge c_*n$. This proves
\[
 q_*(n,k,\eps)\ge c_*\min\{n,k/\eps^2\}
\]
for every parameter triple in RE-03. Together with the $n$-query baseline,
it also proves Corollary~\ref{cor:full}.
